% Generated draft: requires mathematical and source review.
\documentclass[11pt]{article}
\usepackage[margin=25mm]{geometry}
\usepackage{fontspec}
\setmainfont{Latin Modern Roman}
\usepackage{amsmath,amssymb,mathtools,unicode-math}
\setmathfont{Latin Modern Math}
\usepackage{enumitem,needspace,xurl,hyperref}
\setlength{\parindent}{0pt}
\setlength{\parskip}{5pt}
\setlength{\emergencystretch}{3em}
\newcommand{\N}{\mathbb{N}}
\newcommand{\Z}{\mathbb{Z}}
\newcommand{\Q}{\mathbb{Q}}
\newcommand{\R}{\mathbb{R}}
\newcommand{\C}{\mathbb{C}}
\newcommand{\F}{\mathbb{F}}
\newcommand{\A}{\mathbb{A}}
\newcommand{\PP}{\mathbb{P}}
\newcommand{\E}{\mathbb{E}}
\newcommand{\eps}{\varepsilon}
\newcommand{\dd}{\,\mathrm{d}}
\begin{document}
\section*{UnsolvedMath ID: 30004358}
\textbf{Source problem code:} catalogue entry
\subsection{Definitions and mathematical statement}
The supplied record asks for a description of long cycles in Ramsey graphs subject to a restriction on short induced cycles. Its conventions and intended extremal quantity are not fully specified. The following gives a conditional formalization under the usual finite, simple, undirected graph and edge-colouring conventions; these conventions require confirmation.

For a graph \(G\), let \(V(G)\) and \(E(G)\) denote its vertex and edge sets. A cycle of length \(\ell\geq 3\) consists of distinct vertices \(v_1,\ldots,v_\ell\) with edges \(v_iv_{i+1}\) for \(1\leq i<\ell\) and \(v_\ell v_1\). It is induced if the subgraph on these vertices has exactly the cycle edges. For \(r\geq 1\), put \([r]=\{1,\ldots,r\}\). Under the ordinary edge-Ramsey interpretation,
\[
H\to(F)^r
\quad\Longleftrightarrow\quad
\forall c:E(H)\to[r]\;\exists F'\subseteq H\;\exists i\in[r]\quad
\bigl(F'\cong F\ \text{and}\ \forall e\in E(F'),\ c(e)=i\bigr).
\]
Here \(F'\) is a subgraph copy, which need not be induced. If induced Ramsey copies are intended, the requirement \(F'\subseteq H\) must instead require \(F'\) to be an induced subgraph. The supplied notation and title do not resolve this choice.

Interpreting a cycle being ``inherited from a copy of \(F\)'' as containment in a copy of \(F\), the short-cycle condition is
\[
\forall C\subseteq H,\qquad
\bigl(C\text{ is an induced cycle in }H\ \text{and}\ |V(C)|\leq n\bigr)
\ \Longrightarrow\
\exists F_C\subseteq H\quad
\bigl(F_C\cong F\ \text{and}\ C\subseteq F_C\bigr).
\]
The copy \(F_C\) may depend on \(C\). Whether these copies must be induced, must belong to a designated family of copies, or must be the same kind of copies as those in the Ramsey condition is unspecified.

\textbf{Problem as supplied, with unresolved target.} Given a graph \(F\) and integers \(r,n\), describe how many cycles of length greater than \(n\), and which kinds of such cycles, can occur in a graph \(H\) satisfying \(H\to(F)^r\) and the short induced-cycle inheritance condition above. The edge-colouring interpretation requires \(r\geq1\); the source gives no explicit range for \(n\). Under the displayed definition of a cycle, the short-cycle condition is vacuous when \(n<3\). The source does not specify whether the long cycles must themselves be induced, how their number is counted, what constitutes a kind of cycle, or whether the requested result is a universal restriction, an extremal bound, or a construction with controlled long cycles. Thus these formulas describe conditional hypotheses, but the record does not determine a unique precise mathematical target.
\subsection{Short English statement}
For a given graph \(F\) and parameters \(r,n\), the question concerns cycles longer than \(n\) in a graph \(H\) that is Ramsey for \(F\), while every short induced cycle lies inside a copy of \(F\). The intended Ramsey convention, meaning of inheritance, and requested description of long cycles remain unresolved. The supplied partially solved label and dated literature assessment require independent review.
\subsection{Sources}
\begin{itemize}
\item Local UnsolvedMath registry, permanent ID \(30004358\), source problem code \texttt{catalogue entry}, title: \emph{Long Cycles in Induced-Ramsey Graph Constructions}. The record belongs to the Oberwolfach Reports Open Problems collection and records proposed year \(2020\); no individual proposer is supplied.
\item Supplied report citation: \emph{Combinatorics (2021)}, Oberwolfach Reports, DOI \texttt{10.4171/owr/2020/1}, pp.\ 6--89. \url{https://doi.org/10.4171/owr/2020/1}. The background describes Oberwolfach Reports as documenting mathematical workshops and their open-problem discussions. The citation's year and the year appearing in the DOI are retained as supplied.
\item Supplied literature source: Reiher--R\"odl, \emph{Girth Ramsey theory}. \url{https://ems.press/content/serial-article-files/46836?nt=1}.
\item Literature assessment checked at \texttt{2026-08-22}: the supplied curation labels the problem partially solved and states that the girth Ramsey theorem can force every induced cycle up to any prescribed length to lie inside a copy of the target graph, while leaving unrestricted longer-cycle structure unsettled. This is supplied dated curation derived from the OWR collection, not an independently verified literature assessment; its status and cited claims require independent verification.
\item The registry labels the statement \texttt{corrected\_verified}, explaining that mathematical content was preserved while notation, typography, omitted definitions, or duplicate clauses were repaired for self-contained publication. This is a supplied editorial qualification, not independent verification of the present formulation.
\end{itemize}
\end{document}
